\documentclass[11pt]{article}
\usepackage[margin=1in]{geometry}
\usepackage{amsmath,amssymb}
\usepackage[hidelinks]{hyperref}

\newcommand{\APEP}{\operatorname{APEP}}
\newcommand{\Ball}{B}
\newcommand{\norm}[1]{\left\lVert #1 \right\rVert}
\newcommand{\R}{\mathbb R}

\begin{document}

\section*{Human Review: Zero APEPs in Infinite \texorpdfstring{$\ell_1$}{ell-1}-Sums}

\noindent Original automatic proof: \texttt{solution\_packet.pdf}

\noindent Checked date: 17 June 2026

\noindent Status: Manually checked.

\noindent Verdict: The proof appears correct.

\section*{Human Comments}

The proof appears to give a correct negative answer to Question 8.1 of the source paper. The construction takes
\[
  X=\left(\bigoplus_{n=1}^\infty \ell_2^n\right)_1.
\]
Each coordinate space is finite-dimensional, so \(0\) is not an APEP of its unit ball: weak and norm topologies coincide, and a sufficiently small norm-neighbourhood of \(0\) in the ball cannot contain a non-empty slice.

The interesting part is the proof that \(0\) is nevertheless an APEP of \(B_X\). The argument essentially leverages the finite-dimensional structure of the summands together with the fact that a weak neighbourhood of \(0\) in \(B_X\) is controlled by only finitely many functionals. Given such a neighbourhood, determined by \(m\) functionals, one moves to a coordinate \(n>m\). In that coordinate, the \(m\) corresponding coordinate functionals have a non-trivial common kernel. Thus one can choose a unit vector \(u\in \ell_2^n\) invisible to all of them.

The slice is then chosen using the support functional associated to \(u\). Elements in this slice must have almost all of their \(\ell_1\)-sum norm concentrated in the \(n\)-th coordinate. Consequently the other coordinates contribute very little to the finitely many functionals defining the weak neighbourhood. In the \(n\)-th coordinate, the chosen vector \(u\) lies in the common kernel, and the Hilbert-space estimate
\[
  \langle h,u\rangle>1-\delta,\qquad \norm{h}\leq 1=\norm{u}
  \quad\Longrightarrow\quad
  \norm{h-u}\leq \sqrt{2\delta}
\]
shows that the \(n\)-th coordinate is close to \(u\). Hence the contribution from the \(n\)-th coordinate is also small. This proves that the slice lies inside the original weak neighbourhood.

\section*{Extension to Other Finite-Dimensional Blocks}

The proof is not special to Hilbert spaces in its overall structure. The same argument should work with
\[
  X_n=\ell_p^n
\]
for every fixed \(1<p<\infty\). The only replacement needed is for the Hilbert-space estimate above. In \(\ell_p^n\), choose
\[
  u\in S_{\ell_p^n}\cap \bigcap_{j=1}^m \ker a_{j,n}
\]
and then choose a norming functional \(u^*\in S_{(\ell_p^n)^*}\) with \(u^*(u)=1\). Since \(\ell_p^n\) is strictly convex for \(1<p<\infty\), the slice
\[
  \{h\in B_{\ell_p^n}: u^*(h)>1-\delta\}
\]
shrinks to \(u\) as \(\delta\downarrow0\). In finite dimension this can be justified by compactness. Therefore the same estimates go through after replacing the explicit Hilbert bound by this qualitative slice-shrinking fact.

For \(p=1\), the same argument does not apply, since the unit ball has flat faces and a support functional need not expose a single vector. Thus the natural extension of the proof is to the range \(1<p<\infty\).

\section*{Additional Note: Infinite-Dimensional Summands}

If one wants the coordinate spaces themselves to be infinite-dimensional, the same construction should admit a straightforward modification. Take
\[
  X_n=\ell_2^n\oplus_1 c_0
\]
and
\[
  X=\left(\bigoplus_{n=1}^\infty X_n\right)_1.
\]
The \(\ell_2^n\)-part still provides the same finite-dimensional escape direction. Given a weak neighbourhood of \(0\) in \(B_X\) determined by finitely many functionals, choose \(n\) larger than the number of constraints and choose
\[
  u\in S_{\ell_2^n}
\]
orthogonal to the \(\ell_2^n\)-coordinate parts of those functionals. Define the slice using only the \(\ell_2^n\)-component,
\[
  \varphi((h_k,y_k)_k)=\langle h_n,u\rangle.
\]
As before, the condition \(\varphi(x)>1-\delta\) forces almost all of the \(\ell_1\)-sum mass to lie in the \(n\)-th \(\ell_2^n\)-component. It also forces the \(c_0\)-part inside the \(n\)-th coordinate, and all other coordinates, to have small total norm. The same Hilbert-space estimate then puts the slice inside the prescribed weak neighbourhood.

The \(c_0\)-summand is only used to make each \(X_n\) infinite-dimensional while keeping \(0\) visibly non-APEP in each coordinate. One can witness this with a weak neighbourhood in \(B_{X_n}\) which controls both components, for instance
\[
  W_n=\{(h,y)\in B_{\ell_2^n\oplus_1 c_0}:\|h\|_2<1/2,\ |y_1|<1/2\}.
\]
No non-empty slice of \(B_{X_n}\) is contained in \(W_n\): if the \(\ell_2^n\)-part of the defining functional dominates, the slice meets a point of the form \((h,0)\) with \(\|h\|_2=1\); if the \(c_0\)-part dominates, the slice meets a point of the form \((0,y)\) with \(|y_1|=1\), using the standard \(\ell_1\)-description of \(c_0^*\). Thus \(0\) is not an APEP of \(B_{X_n}\) for any \(n\), while the \(\ell_1\)-sum still has \(0\) as an APEP.

\section*{Review Outcome}

Archive this packet as an apparently correct counterexample. The proof's central mechanism is sound: every coordinate fails APEP at \(0\) because it is finite-dimensional, while the infinite \(\ell_1\)-sum allows one to escape any finite list of weak constraints by moving to a sufficiently high-dimensional coordinate and choosing a direction annihilated by those constraints. The same construction should also work with \(\ell_p^n\) blocks for \(1<p<\infty\).

\end{document}
